\documentclass{exam}[12pt]

\usepackage{amsfonts,amsmath,amsthm,amssymb}


\boxedpoints

\addpoints
\pointsinmargin

\begin{document}
\pagestyle{head}

\firstpageheader{\large Calculus II}{\bf\Large Exam I }{\large Name:\hspace{1in} }
\small{\textit{Calculators are not allowed on this  exam. 
Work all questions completely.  Show all work as described in class. \\
Copyright 2011 Dr. Brock Williams and
Texas Tech University.  Unauthorized reproduction prohibited.\\Point values for each
problem are in the boxes in the margin.}}

\begin{questions}
\large{

\question[7] Set up an integral to find the area of the region bounded by
$y=\sqrt{x}$ and $y= \frac{1}{3} x$.
\vspace{1.3in}

\question Set up an integral to find the volume of the solid generated when
the region bounded by $y=\sqrt{x}$ and $y= \frac{1}{3} x$ is
\begin{parts}
\part[7] Rotated about the $x$-axis using washers
\vspace{1.7in}
\part[7] Rotated about the $x$-axis using shells
\vspace{1.7in}
\part[7] Rotated about the line $x=-4$
\vspace{1.7in}
\part[7] Rotated about the line $y=5$
\vspace{1.7in}
%\part[7] Rotated about the line $x= -4$
%\vspace{1in}
\end{parts}

\question[14] Set up integrals to find the mass, the moment about the $x$-axis, and the
moment about the $y$-axis of the region bounded by $y=\sqrt{x}$ and $y= \frac{1}{3} x$.
Give the formula to find the coordinates of the center of mass. (Be sure to indicate which 
integral is which.)
\vspace{4in}

\question[2] \textbf{BONUS.} Set up an integral to find the moment about the line $x=-3$ of the
region bounded by $y=\sqrt{x}$ and $y= \frac{1}{3} x$.  
\vspace{1.0in}

\question[10] Graph the function $r=3(1-\cos\theta),$ $0\leq\theta\leq\2pi$ 
and set up an
integral to find the area enclosed by the graph.
\vspace{2.2in}

\pagebreak

\question[8] Consider the curve given by $f(x)=\displaystyle \sqrt{x}, \quad 0\leq x\leq 3$.
Set up an integral to find the length of this curve.
\vspace{1.8in}


\question[6] A spring is naturally 5 ft long and has spring constant $k=70$ lb/ft.  Set up an 
integral for the work done in stretching the spring from 6 ft long to 8 ft long.
\vspace{1.5in}


\question[12] A conical tank has a height of 8 feet and a radius of 7 feet.
It is filled with a fluid having density 63 lb/ft$^3$.
Set up an integral to find the work done in pumping the fluid out of the tank if
the tank is filled to a depth of 5 feet and the water must also be
 pumped out a 4 foot pipe sticking
out the top of the tank.



%\question[10] A dam has a rectangular gate with dimensions as shown below.
%The top of the gate is 2 feet below the water level.  Set up an integral
%to find the force on
%the gate. (Recall the density of water is 62.4 lb/ft$^3$.)

%\question[10] A dam is 20 ft tall and is in the shape of the region bounded by 
%the parabola $y=\displaystyle \frac{1}{5} x^2$ and the line $y=20$ as shown below.
%The water is level with the top of the dam.  Set up an integral
%to find the force of the water on
%the dam. (Recall the density of water is 62.4~lb/ft$^3$.)

\pagebreak
\question  
If the amount of happiness while working a math test is constant,
then the total joy produced during the test is given by
\[
\text{Joy = $\left(\sqrt{{\text{Happiness}}}\right)$ * (The Time The Happiness Lasts)}.
\]
In a certain 50 minute Calculus test, the Happiness was given by $\displaystyle H(t) = 1 - \frac{t}{50}$, where
$0\leq t\leq 50$.
\begin{parts}
\part[10] Describe in complete detail the fundamental idea of Chapter 6 and explain
 in complete detail how this fundamental idea can be
applied to find the total Joy in the test.
\vspace{6.5in}
\part[5] Set up an integral to find the total Joy in the test.
\end{parts}

%\question If the ammount of mayhem is constant in an episode of the
%Muppet Show, there is a simple formula that gives the amount of mess created:
%\[
%\text{Mess = Mayhem * Time}
%\]
%That is, the mess is the product of the mayhem and the length of time it lasts.
%Typically, however, the mayhem varies over time.  

}
\end{questions}
\end{document}



